\documentclass[11pt,a4paper]{article}
\usepackage[T1]{fontenc}
\usepackage{lmodern}
\usepackage[margin=27mm]{geometry}
\usepackage{amsmath,amssymb,amsthm,mathtools}
\usepackage{microtype,booktabs,array}
\usepackage{xcolor}
\definecolor{linkblue}{RGB}{31,67,104}
\usepackage[colorlinks=true,linkcolor=linkblue,citecolor=linkblue,urlcolor=linkblue]{hyperref}
\usepackage{fancyhdr}
\pagestyle{fancy}\fancyhf{}
\fancyhead[L]{\small Le Roy functions: positivity beyond infinite divisibility}
\fancyhead[R]{\small Research draft}
\fancyfoot[C]{\thepage}
\setlength{\headheight}{14pt}
\setlength{\emergencystretch}{2em}
\numberwithin{equation}{section}
\newtheorem{theorem}{Theorem}[section]
\newtheorem{proposition}[theorem]{Proposition}
\newtheorem{lemma}[theorem]{Lemma}
\newtheorem{corollary}[theorem]{Corollary}
\theoremstyle{definition}\newtheorem{definition}[theorem]{Definition}
\theoremstyle{remark}\newtheorem{remark}[theorem]{Remark}
\newcommand{\R}{\mathbb R}
\newcommand{\C}{\mathbb C}
\newcommand{\E}{\mathbb E}
\newcommand{\ii}{\mathrm i}
\newcommand{\dd}{\,\mathrm d}
\DeclareMathOperator{\Rea}{Re}
\DeclareMathOperator{\Ima}{Im}
\DeclareMathOperator{\supp}{supp}
\title{Complete monotonicity beyond logarithmic infinite divisibility\\[5pt]
\large A counterexample theorem for Le Roy functions}
\author{Keio Research Notes\thanks{AI-assisted research manuscript prepared with OpenAI Codex. The argument received separate internal mathematical checks; these are not external peer review. Historical priority is not certified.}}
\date{3 November 2026}
\hypersetup{pdftitle={Complete monotonicity beyond logarithmic infinite divisibility},pdfauthor={Keio Research Notes},pdfsubject={A counterexample theorem for Simon's Le Roy complete-monotonicity conjecture}}

\begin{document}
\maketitle
\begin{abstract}
For the Le Roy function
\(F_{\alpha,\beta}^{(\gamma)}(z)=\sum_{n\ge0}z^n/\Gamma(\beta+\alpha n)^\gamma\),
we prove a strict separation between complete monotonicity on the negative
half-line and logarithmic infinite divisibility of its representing law.
For every fixed real \(\gamma>1\) and \(0<\alpha<1/\gamma\), let
\(\beta_0\) be the threshold at which
\(z+\gamma(z^{\beta-\alpha}-z^\beta)\le1\) holds for all \(0<z<1\).
There is a nonempty interval immediately below \(\beta_0\) on which the
Le Roy function remains completely monotone. This disproves the universal
equivalence proposed by Thomas Simon. The proof combines strict positivity
of an infinitely divisible density, obtained from an analytic convolution
root, with uniform estimates at both endpoints for an inverse Mellin
transform. These estimates allow positivity to persist as the signed
L\'evy kernel becomes negative. The argument covers nonintegral
\(\gamma\), is independent of numerical computation, and concerns one
previously posed conjecture. It does not determine the exact complete
monotonicity region.
\end{abstract}

\noindent\textbf{Keywords.} Complete monotonicity; Le Roy function; Mellin
inversion; infinite divisibility; positive densities; gamma moments.

\section{The original conjecture and the result}

A smooth real function \(h\) on \((0,\infty)\) is completely monotone
(CM) if \((-1)^n h^{(n)}(x)\ge0\) for all integers \(n\ge0\).
A positive random variable is logarithmically infinitely divisible
(log-ID) if its logarithm has an infinitely divisible law on \(\R\).
Consider the entire function
\begin{equation}\label{eq:F}
 F_{\alpha,\beta}^{(\gamma)}(w)
 =\sum_{n=0}^{\infty}\frac{w^n}{\Gamma(\beta+\alpha n)^\gamma},
 \qquad \alpha,\beta,\gamma>0.
\end{equation}
Simon \cite[Section 3]{Simon} proposed the equivalence
\begin{equation}\label{eq:original}
 F_{\alpha,\beta}^{(\gamma)}(-x)\text{ is CM}
 \quad\Longleftrightarrow\quad
 z+\gamma(z^{\beta-\alpha}-z^\beta)\le1
 \quad\text{for every }0<z<1.
\end{equation}
The inequality is sufficient by a L\'evy--Khintchine construction.
Its proposed necessity is the target here. The original parameters,
function, domain, and universal quantifier in \eqref{eq:original} are
retained.

For fixed \(\gamma>1\), \(0<\alpha<1/\gamma\), define
\begin{equation}\label{eq:B}
 B(t)=\alpha+\frac{\alpha}{t}
 \log\!\left(\frac{\gamma(1-e^{-t})}{1-e^{-t/\alpha}}\right),
 \qquad
 \beta_0=\max_{t>0}B(t).
\end{equation}
The existence of this maximum and its relationship to
\eqref{eq:original} are proved below.

\begin{theorem}[Strict separation]\label{thm:main}
Fix any real \(\gamma>1\) and \(0<\alpha<1/\gamma\). Then
\[
 \alpha<\beta_0<\frac{\alpha(1+\gamma)}2.
\]
There exists \(\varepsilon>0\), with \(\beta_0-\varepsilon>\alpha\),
such that for every \(\beta\in(\beta_0-\varepsilon,\beta_0)\):
\begin{enumerate}
\item \(F_{\alpha,\beta}^{(\gamma)}(-x)\) is completely monotone on
\((0,\infty)\);
\item the inequality on the right of \eqref{eq:original} fails for some
\(z\in(0,1)\);
\item there is a positive random variable \(X_\beta\), not log-ID, with
\begin{equation}\label{eq:moment}
 \E X_\beta^s
 =\frac{\Gamma(1+s)\Gamma(\beta)^\gamma}
 {\Gamma(\beta+\alpha s)^\gamma},\qquad \Rea s>-1,
\end{equation}
and \(\E e^{wX_\beta}=\Gamma(\beta)^\gamma
F_{\alpha,\beta}^{(\gamma)}(w)\) for every \(w\in\C\).
\end{enumerate}
\end{theorem}

The theorem gives a negative resolution of the literal equivalence
\eqref{eq:original}. The interval width is established to be positive;
no explicit numerical width or particular rational value of \(\beta\)
is asserted. Infinitely many parameter choices remain one conjecture
resolution, not multiple independent targets.

The transform representation is classical in the rational/integer
specializations studied by G\'orska, Horzela, and Garrappa \cite{GHG}.
The argument here requires uniformity in \(\beta\), strict positivity at
the log-ID threshold, and continuation across that threshold. These are
proved directly, including for noninteger \(\gamma\).

\section{The threshold and the signed L\'evy kernel}

Fix \(\alpha,\gamma\) as in Theorem~\ref{thm:main}, and write
\(\nu=1-\alpha\gamma\in(0,1)\). For \(\beta>\alpha\), put
\begin{equation}\label{eq:phi}
 \phi_\beta(t)=\frac{e^{-t}}{1-e^{-t}}
 -\gamma\frac{e^{-\beta t/\alpha}}{1-e^{-t/\alpha}},\qquad t>0.
\end{equation}
Elementary rearrangement gives
\begin{equation}\label{eq:phiB}
 \phi_\beta(t)\ge0\quad\Longleftrightarrow\quad\beta\ge B(t).
\end{equation}
With \(z=e^{-t/\alpha}\), the same inequality becomes precisely the
right side of \eqref{eq:original}.

\begin{lemma}\label{lem:threshold}
The maximum in \eqref{eq:B} is attained in \((0,\infty)\), and
\(\alpha<\beta_0<\alpha(1+\gamma)/2\). The scalar inequality in
\eqref{eq:original} holds for all \(z\in(0,1)\) exactly when
\(\beta\ge\beta_0\).
\end{lemma}
\begin{proof}
As \(t\downarrow0\), the logarithm in \eqref{eq:B} tends to
\(\log(\alpha\gamma)<0\), so \(B(t)\to-\infty\). As
\(t\to\infty\),
\[
 B(t)=\alpha+\frac{\alpha\log\gamma}{t}+o(t^{-1}).
\]
Thus \(B\) is greater than \(\alpha\) somewhere and tends to
\(\alpha\) at infinity. Continuity implies that its maximum is
attained at a finite positive point and exceeds \(\alpha\).

Let \(\beta_U=\alpha(1+\gamma)/2\). Since \(1/\alpha>\gamma\),
\begin{align*}
 \phi_{\beta_U}(t)
 &>\frac{e^{-t}}{1-e^{-t}}
 -\gamma\frac{e^{-(1+\gamma)t/2}}{1-e^{-\gamma t}}\\
 &=\frac{e^{-t/2}}2\left\{
 \frac1{\sinh(t/2)}-\frac\gamma{\sinh(\gamma t/2)}\right\}>0.
\end{align*}
The last inequality follows because \(u\mapsto\sinh(u)/u\) is
strictly increasing for \(u>0\). Hence \(B(t)<\beta_U\) at every
positive \(t\). Since its maximum is attained, \(\beta_0<\beta_U\).
The equivalence follows from \eqref{eq:phiB}.
\end{proof}

For complex \(s\) with \(\Rea s>-1\), define
\begin{equation}\label{eq:M}
 M_\beta(s)=\Gamma(1+s)\Gamma(\beta)^\gamma
 \exp\{-\gamma\operatorname{Log}\Gamma(\beta+\alpha s)\}.
\end{equation}
Here \(\operatorname{Log}\Gamma\) is the holomorphic logarithm on the
right half-plane that is real on the positive real axis. This convention
is essential when \(\gamma\notin\mathbb N\). It is available because
\(\Rea(\beta+\alpha s)>\beta-\alpha>0\).

The logarithmic gamma integral gives
\begin{equation}\label{eq:LK}
 \log M_\beta(s)
 =d_\beta s+\int_0^\infty
 (e^{-st}-1+st)\phi_\beta(t)\,\frac{\mathrm dt}{t},
 \qquad d_\beta=\psi(1)-\alpha\gamma\psi(\beta).
\end{equation}
For clarity, this identity can be checked without choosing an implicit
complex power. Twice differentiating both sides gives
\(\psi'(1+s)-\alpha^2\gamma\psi'(\beta+\alpha s)\), using the
standard integral for the trigamma function. The two sides agree at
\(s=0\) and have the same derivative there. Analyticity proves the
identity throughout \(\Rea s>-1\).

At the endpoints,
\begin{equation}\label{eq:phitails}
 \phi_\beta(t)=\frac\nu t+O(1)\quad(t\downarrow0),\qquad
 \phi_\beta(t)=O(e^{-t})\quad(t\to\infty).
\end{equation}
Consequently \(\phi_\beta(t)\,\mathrm dt/t\) is a signed L\'evy
measure with finite absolute second moment. At \(\beta=\beta_0\)
it is nonnegative, so \eqref{eq:LK} with \(s=\ii u\) is the
characteristic exponent of an infinitely divisible law. Its jumps
correspond to \(-t\), and the compensation \(+\ii ut\) in
\eqref{eq:LK} is admissible by \eqref{eq:phitails}.

\section{An auxiliary positivity-continuation problem}

The useful auxiliary question is the following: when can a family of
inverse transforms remain nonnegative after its logarithmic transform
loses a nonnegative L\'evy representation? Two elementary facts give a
criterion. They are stated explicitly to separate the general argument
from its special-function estimates.

\begin{lemma}[Strict positivity from a convolution root]\label{lem:analyticID}
Let \(\chi\) be an infinitely divisible characteristic function on
\(\R\). Suppose
\[
 |\chi(u)|\le C(1+|u|)^A e^{-a|u|}\quad(u\in\R)
\]
for some \(a,C>0\) and finite \(A\). Its probability density is
continuous and strictly positive at every real point.
\end{lemma}
\begin{proof}
Infinite divisibility supplies a characteristic function \(\theta\)
with \(\theta^2=\chi\), and
\[
 |\theta(u)|\le C^{1/2}(1+|u|)^{A/2}e^{-a|u|/2}.
\]
Fourier inversion gives a bounded continuous nonnegative probability
density \(p\) for this convolution root. The same integral defines a
holomorphic function in \(|\Ima w|<a/2\). Thus \(p\) is real
analytic and not identically zero. Its real zeros are isolated, so
\(p(v)p(w-v)>0\) outside a discrete set for every fixed real \(w\).
It follows that
\[
 (p*p)(w)=\int_\R p(v)p(w-v)\dd v>0.
\]
This integral is finite since \(p\) is bounded and integrable. The
convolution is the continuous density with characteristic function
\(\chi\).
\end{proof}

\begin{lemma}[Continuation with uniform endpoint control]\label{lem:continue}
Let \(q_\theta(x)\) be jointly continuous for \(x>0\) and real
\(\theta\) in a neighborhood of \(\theta_0\). Suppose
\(q_{\theta_0}(x)>0\) for every \(x>0\). If there are
\(0<x_0<x_1<\infty\) such that
\[
 q_\theta(x)>0\quad\text{on }(0,x_0]\cup[x_1,\infty)
\]
uniformly for \(\theta\) in that neighborhood, then
\(q_\theta(x)>0\) for all \(x>0\) whenever
\(|\theta-\theta_0|\) is sufficiently small.
\end{lemma}
\begin{proof}
The continuous function \(q_{\theta_0}\) has a positive minimum on
\([x_0,x_1]\). Uniform continuity on a compact parameter neighborhood
times this interval preserves, for nearby \(\theta\), at least half
that minimum. The assumed endpoint control covers the rest of the
positive axis.
\end{proof}

This compactness argument alone does not prove the theorem: pointwise
positivity or pointwise endpoint asymptotics would not supply its
uniform hypothesis. The next section establishes precisely that
hypothesis for \eqref{eq:M}.

\section{Inverse Mellin densities and uniform tails}

Let \(K\) be any compact interval in \((\alpha,\infty)\).
All constants and uniform remainders in this section may depend on
\(\alpha,\gamma,K\). Define the real-valued candidate density
\begin{equation}\label{eq:q}
 q_\beta(x)=\frac1{2\pi\ii x}
 \int_{-\ii\infty}^{\ii\infty}M_\beta(s)x^{-s}\dd s,
 \qquad x>0,\quad\beta\in K.
\end{equation}
It need not initially be nonnegative. Stirling's formula on each fixed
vertical line supplies a polynomial bound times
\(e^{-\pi\nu|\Ima s|/2}\), uniformly for \(\beta\in K\).
The integral therefore converges absolutely and is jointly continuous
in \((\beta,x)\). Reality follows by complex conjugation.

\subsection{The endpoint at zero}

Choose
\begin{equation}\label{eq:eta}
 0<\eta<\min\{1,(\min K)/\alpha-1\}.
\end{equation}
On the whole strip \(\Rea s\ge-1-\eta\), the denominator in
\eqref{eq:M} uses a single holomorphic logarithm, because
\(\Rea(\beta+\alpha s)>0\). Shifting the contour in
\eqref{eq:q} to \(\Rea s=-1-\eta\) crosses only the simple
numerator pole at \(s=-1\). The horizontal integrals tend to zero
by the vertical Stirling bounds. Taking the residue yields
\begin{equation}\label{eq:small}
 q_\beta(x)=A_\beta+O_K(x^\eta),\qquad
 A_\beta=\frac{\Gamma(\beta)^\gamma}
 {\Gamma(\beta-\alpha)^\gamma}>0,
 \qquad x\downarrow0.
\end{equation}
The remainder is uniform because the shifted integral is bounded
absolutely and uniformly after factoring out \(x^\eta\).
In particular, \(\min_{\beta\in K}A_\beta>0\).

\subsection{The endpoint at infinity}

Introduce
\begin{align}
 \kappa_\beta&=\frac{1+\gamma}{2}-\gamma\beta,&
 C_\beta&=(2\pi)^{(1-\gamma)/2}\Gamma(\beta)^\gamma
 \alpha^{\gamma(1/2-\beta)},\label{eq:constants}\\
 r&=(\alpha^{\alpha\gamma}x)^{1/\nu}.&&\label{eq:r}
\end{align}
Uniform sectorial Stirling estimates \cite[Section 5.11]{DLMF} give,
as \(|s|\to\infty\) in the right half-plane,
\begin{equation}\label{eq:stirling}
 M_\beta(s)=C_\beta s^{\kappa_\beta}
 \exp\{\nu s(\log s-1)-\alpha\gamma s\log\alpha\}
 \bigl(1+O_K(|s|^{-1})\bigr).
\end{equation}
Here \(\log s\) is the principal logarithm in that half-plane. The
shifted gamma arguments vary in a compact family, so the remainder is
uniform in \(\beta\). Taking the fixed real power \(\gamma\)
through the holomorphic logarithm preserves this estimate.

\begin{proposition}[Uniform saddle estimate]\label{prop:large}
Uniformly for \(\beta\in K\),
\begin{equation}\label{eq:large}
 q_\beta(x)=\frac{C_\beta}{x\sqrt{2\pi\nu}}
 r^{\kappa_\beta+1/2}e^{-\nu r}\bigl(1+o_K(1)\bigr),
 \qquad x\to\infty.
\end{equation}
\end{proposition}
\begin{proof}
Move the contour in \eqref{eq:q} to \(\Rea s=r\). No poles
are crossed; for each fixed \(r\), horizontal integrals vanish as
their imaginary height tends to infinity. Because
\(\log x+\alpha\gamma\log\alpha=\nu\log r\), the exponential
part of \(M_\beta(s)x^{-s}\) is
\[
 \exp\{\nu s(\log(s/r)-1)\}.
\]
Set \(s=r+\ii\sqrt r\,u\). Expansion around \(s=r\) gives
\begin{equation}\label{eq:saddlelimit}
 \frac{M_\beta(r+\ii\sqrt r\,u)x^{-r-\ii\sqrt r\,u}}
 {C_\beta r^{\kappa_\beta}e^{-\nu r}}
 \longrightarrow e^{-\nu u^2/2}
\end{equation}
for each fixed \(u\), uniformly for \(\beta\in K\). In detail,
the first derivative of the exponent vanishes, its second derivative
is \(\nu/r\), and multiplication by \((\ii\sqrt r\,u)^2/2\)
produces \(-\nu u^2/2\).

To control the entire contour, let \(v=u/\sqrt r\) and define
\[
 H(v)=\tfrac12\log(1+v^2)-v\arctan v.
\]
The modulus of the left side of \eqref{eq:saddlelimit} is bounded by
\begin{equation}\label{eq:dom}
 D_K(1+v^2)^{\kappa_\beta/2}e^{\nu r H(v)}.
\end{equation}
The function \(H\) is even, \(H(0)=0\), and
\(H'(v)=-\arctan v\) for \(v>0\). Hence
\begin{align}
 H(v)&\le-v^2/4&&(|v|\le1),\label{eq:Hsmall}\\
 H(v)&\le H(1)-\tfrac\pi4(|v|-1)&&(|v|\ge1),\label{eq:Hlarge}
\end{align}
with \(H(1)<0\). The factor involving \(\kappa_\beta\) is
uniformly bounded on \(|v|\le1\); this part of the integrand is
dominated by a constant times \(e^{-\nu u^2/4}\).

For the outer region put
\(\kappa_+=\max(0,\max_{\beta\in K}\kappa_\beta)\).
Since
\(\log(1+v^2)\le\log2+(|v|-1)\) for \(|v|\ge1\),
\eqref{eq:dom}--\eqref{eq:Hlarge}, for all sufficiently large
\(r\), are bounded by
\[
 D'_K e^{\nu rH(1)}
 \exp\{-\nu\pi r(|v|-1)/8\}.
\]
Its integral over \(|u|\ge\sqrt r\) is
\(O_K(r^{-1/2}e^{\nu rH(1)})\). Thus the outer contour is
uniformly negligible. Dominated convergence on the central region,
uniform on bounded \(u\)-intervals in \(\beta\), now gives the
uniform integral limit
\(\int_\R e^{-\nu u^2/2}\dd u=\sqrt{2\pi/\nu}\).
Finally \(\mathrm ds=\ii\sqrt r\,\mathrm du\) and the prefactor
\(1/(2\pi\ii x)\) in \eqref{eq:q} yield
\eqref{eq:large}.
\end{proof}

All leading factors in \eqref{eq:small} and \eqref{eq:large} are
positive. Therefore, for this entire compact interval \(K\), there
are \(0<x_0<x_1<\infty\) such that
\begin{equation}\label{eq:endpositive}
 q_\beta(x)>0\quad
 (\beta\in K,\ x\in(0,x_0]\cup[x_1,\infty)).
\end{equation}
This is a uniform conclusion on an unbounded domain, not an inference
from a finite numerical grid.

\subsection{Normalization before positivity}

The endpoint estimates imply absolute integrability of
\(x^s q_\beta(x)\) for every \(\Rea s>-1\). Let
\(p_\beta(y)=e^y q_\beta(e^y)\). The definition
\eqref{eq:q} says that \(p_\beta\) is the inverse Fourier
transform of the integrable function \(u\mapsto M_\beta(\ii u)\).
The endpoint bounds also give \(p_\beta\in L^1(\R)\). Fourier
inversion therefore identifies its Fourier transform with
\(M_\beta(\ii u)\). Both sides of
\begin{equation}\label{eq:mellin}
 \int_0^\infty x^s q_\beta(x)\dd x=M_\beta(s)
\end{equation}
are holomorphic on \(\Rea s>-1\), so the identity theorem extends
this equality from the imaginary axis to that half-plane. In particular,
\begin{equation}\label{eq:mass}
 \int_0^\infty q_\beta(x)\dd x=1.
\end{equation}
No nonnegativity assumption was used in this normalization.

\section{Proof of strict separation}

Choose a compact interval \(K\subset(\alpha,\infty)\) with
\(\beta_0\) in its interior. By Section~2,
\(M_{\beta_0}(\ii u)\) is an infinitely divisible characteristic
function. Vertical Stirling estimates give
\[
 |M_{\beta_0}(\ii u)|
 \le D(1+|u|)^A e^{-\pi\nu|u|/2}
\]
for some finite \(A\) and \(D>0\). Lemma~\ref{lem:analyticID}
shows that its density is strictly positive on \(\R\). Fourier
uniqueness identifies that density with
\(e^y q_{\beta_0}(e^y)\), so
\begin{equation}\label{eq:strict}
 q_{\beta_0}(x)>0\qquad(x>0).
\end{equation}
This step uses infinite divisibility only at the threshold itself.

Joint continuity, \eqref{eq:endpositive}, and
Lemma~\ref{lem:continue} imply that \(q_\beta\) remains strictly
positive everywhere for \(|\beta-\beta_0|<\varepsilon\), with
some \(\varepsilon>0\) small enough that this interval lies in
\(K\). By \eqref{eq:mass}, these functions are probability
densities. Let \(X_\beta\) have density \(q_\beta\).
Equation~\eqref{eq:mellin} proves \eqref{eq:moment}.

Because \(0<\nu<1\), the stretched exponential in
\eqref{eq:large} decays faster than \(e^{-cx}\) for every fixed
\(c>0\). Thus \(\E e^{cX_\beta}<\infty\) for all \(c>0\),
and the exponential series can be integrated absolutely for every
\(w\in\C\). It follows that
\begin{align}
 \E e^{wX_\beta}
 &=\sum_{n\ge0}\frac{w^n}{n!}M_\beta(n)
 =\Gamma(\beta)^\gamma F_{\alpha,\beta}^{(\gamma)}(w),
 \label{eq:mgf}\\
 F_{\alpha,\beta}^{(\gamma)}(-x)
 &=\Gamma(\beta)^{-\gamma}\E e^{-xX_\beta}.
 \label{eq:laplace}
\end{align}
Differentiating the Laplace integral proves complete monotonicity.
For \(\beta<\beta_0\), Lemma~\ref{lem:threshold} proves failure
of the unchanged scalar inequality in \eqref{eq:original}.

It remains to justify explicitly the stronger non-log-ID conclusion.
For such a \(\beta\), continuity and \eqref{eq:phiB} show that
\(\phi_\beta\) is negative on a nonempty open interval.
Let \(V=\log X_\beta\). Its characteristic function is
\(M_\beta(\ii u)\), and \eqref{eq:LK} is its continuous
characteristic logarithm, equal to zero at \(u=0\). Twice
differentiating yields
\begin{equation}\label{eq:fouriermeasure}
 -\frac{\mathrm d^2}{\mathrm du^2}\log M_\beta(\ii u)
 =\int_0^\infty e^{-\ii ut}\,t\phi_\beta(t)\dd t.
\end{equation}
The signed measure on the right has finite total variation by
\eqref{eq:phitails}. Also \(\E V^2<\infty\), directly from
\eqref{eq:small} and \eqref{eq:large}.

If \(V\) were infinitely divisible, its finite variance would give a
Gaussian coefficient \(\sigma^2\ge0\) and a nonnegative L\'evy
measure \(\Pi\) with \(\int y^2\Pi(\mathrm dy)<\infty\).
The left side of \eqref{eq:fouriermeasure} would then be the Fourier
transform of the finite nonnegative measure
\[
 \sigma^2\delta_0(\mathrm dy)+y^2\Pi(\mathrm dy).
\]
Uniqueness of Fourier transforms of finite measures identifies this
measure with the pushforward of \(t\phi_\beta(t)\,\mathrm dt\)
under \(t\mapsto-t\). The latter is negative on an interval, a
contradiction. This completes the proof of
Theorem~\ref{thm:main}. \qed

\section{A concrete specialization and consistency checks}

For \(\alpha=1/3\) and \(\gamma=2\), setting \(z=t^3\) in
the scalar inequality and dividing by \(1-t\) reduces it to
\[
 2t^{3\beta-1}\le1+t+t^2\qquad(0<t<1).
\]
Thus
\begin{equation}\label{eq:bexplicit}
 b_*:=\beta_0(1/3,2)
 =\frac13\left(1+
 \max_{0<t<1}\frac{\log((1+t+t^2)/2)}{\log t}\right).
\end{equation}
The quotient tends to zero from above as \(t\downarrow0\) and
to \(-\infty\) as \(t\uparrow1\). Its maximum is attained.
Furthermore
\[
 1+t+t^2-2\sqrt t=(1-\sqrt t)^2+t^2>0
\]
implies \(1/3<b_*<1/2\). Numerical evaluation gives
\(b_*\approx0.4347490386\); this decimal is illustrative and is not
used to certify any particular parameter as a counterexample.

For an independent check of the constants, let
\(Z=X^3/27\). The gamma triplication formula gives the inverse
Mellin density
\begin{equation}\label{eq:G}
 f_\beta(z)=\frac{\Gamma(\beta)^2}{\Gamma(1/3)\Gamma(2/3)}
 G^{3,0}_{2,3}\!\left(z\;\middle|\;
 \begin{matrix}\beta-1,\beta-1\\[-2pt]-2/3,-1/3,0\end{matrix}\right).
\end{equation}
Here the Meijer function may simply be defined by the absolutely
convergent contour integral
\[
 \frac1{2\pi\ii}\int_{1-\ii\infty}^{1+\ii\infty}
 \frac{\Gamma(s-2/3)\Gamma(s-1/3)\Gamma(s)}
 {\Gamma(s+\beta-1)^2}z^{-s}\dd s.
\]
The transformation of densities is
\(f_\beta(z)=z^{-2/3}q_\beta(3z^{1/3})\). Accordingly,
\eqref{eq:small} and \eqref{eq:large} give
\begin{align}
 f_\beta(z)&=
 \frac{\Gamma(\beta)^2}{\Gamma(\beta-1/3)^2}z^{-2/3}
 (1+o_K(1)),&&z\downarrow0,\label{eq:Gsmall}\\
 f_\beta(z)&=
 \frac{\Gamma(\beta)^2}{\Gamma(1/3)\Gamma(2/3)}
 z^{1-2\beta}e^{-z}(1+o_K(1)),&&z\to\infty.
 \label{eq:Glarge}
\end{align}
These agree with the special contour calculation. The general proof,
however, needs neither rational \(\alpha\) nor integer \(\gamma\).

As a further normalization check on the saddle constants, the formulas
with \(\alpha=1/2,\gamma=1,\beta=1\) give
\(q(x)=\pi^{-1/2}e^{-x^2/4}\), the known half-normal density.
This check is outside the counterexample regime and is not an
additional conjecture resolution.

\section{Prior work, the corpus, and the scope of the claim}

\paragraph{Classical inputs and the proposed contribution.}
Mellin and Fourier inversion, gamma asymptotics, L\'evy--Khintchine
theory, and compactness are classical. G\'orska, Horzela, and Garrappa
\cite{GHG} obtained inverse Meijer representations and pointwise
asymptotics in rational/integer cases. Their numerical positivity
thresholds do not provide the uniform continuation proof above.
The contribution proposed here is the strict gap theorem, established
by the analytic convolution-root argument together with uniform
endpoint control for all real parameters in the stated regime.

Gamma-type probability distributions with non-ID logarithms are not
new in general. Dufresne \cite[Example 2]{Dufresne} gives such an
example, and Kadankova, Simon, and Wang \cite{KSW} construct related
signed spectral measures. Wang \cite{Wang} studies further gamma
moment and quasi-infinite-divisibility questions. The theorem here
concerns the specific family and universal equivalence
\eqref{eq:original}; it does not claim the first gamma-type non-ID law.

\paragraph{How the supplied corpus was used.}
The Verifiable Mathematical Reasoning Corpus \cite{Corpus} was
examined at commit
\texttt{ecd87de3348dcb706a76bb12fde722d7a348a012}.
Manuscript [067], especially its method M9, suggests passing from
derivative signs to Laplace-density signs. Manuscript [096] emphasizes
exact integral representations, endpoint analysis, and separation of
classical inputs from a new application. Those ideas motivated the
auxiliary problem in Section~3. No unverified theorem from the corpus
is a premise of the proof; the required estimates and positivity
arguments are given in full here.

\paragraph{Priority and verification.}
A targeted literature search through 3 November 2026 examined the
original paper, its relevant gamma-moment references, the earlier
Meijer representations, later gamma-moment work, and exact-title and
conjecture searches. It did not locate an identical resolution of
\eqref{eq:original}. This is a report about the search performed, not
a proof of the absence of prior work. Historical priority and external
peer validation remain unconfirmed. The separate question concerning
convexity of Simon's scalar threshold is not claimed as a result here.

The proof was checked in separate internal derivations of the
\((1/3,2)\) case and the full real-parameter case. The checks covered
branch choices, residues, uniform saddle bounds, normalization,
strict positivity, and signed L\'evy uniqueness. Numerical samples
were exploratory only. No finite search, assumed positivity theorem,
or unproved conjecture is needed in the proof.

\paragraph{Research count and remaining work.}
This manuscript contains a complete analytic argument against one
original conjecture. The broader request for ten difficult conjecture
resolutions has not been achieved. Other investigated problems in
number theory, graph theory, and symbolic integration produced no
additional unrestricted resolution. Restricted classes, already
published solutions, and multiple instances of the same family have
not been counted. The exact CM threshold for Le Roy functions remains
undetermined by this paper, as does a certified numerical width of
the counterexample interval.

\begin{thebibliography}{9}

\bibitem{Simon}
T. Simon, \emph{Remark on a Mittag--Leffler function of Le Roy type},
Integral Transforms and Special Functions \textbf{33} (2022), no.~2,
108--114.
\href{https://doi.org/10.1080/10652469.2021.1913138}{doi:10.1080/10652469.2021.1913138}.
Author version: \href{https://arxiv.org/abs/2012.11018v2}{arXiv:2012.11018v2},
25 March 2021; conjecture in Section~3.

\bibitem{GHG}
K. G\'orska, A. Horzela, and R. Garrappa,
\emph{Some results on the complete monotonicity of the Mittag--Leffler
functions of Le Roy type}, Fractional Calculus and Applied Analysis
\textbf{22} (2019), no.~5, 1284--1306.
\href{https://doi.org/10.1515/fca-2019-0068}{doi:10.1515/fca-2019-0068}.
\href{https://arxiv.org/abs/1912.07100v2}{arXiv:1912.07100v2}.

\bibitem{Dufresne}
D. Dufresne, \emph{G distributions and the beta-gamma algebra},
Electronic Journal of Probability \textbf{15} (2010), no.~71,
2163--2199.
\href{https://doi.org/10.1214/EJP.v15-845}{doi:10.1214/EJP.v15-845}.

\bibitem{KSW}
T. Kadankova, T. Simon, and M. Wang,
\emph{On some new moments of Gamma type},
Statistics \& Probability Letters \textbf{165} (2020), 108854.
\href{https://doi.org/10.1016/j.spl.2020.108854}{doi:10.1016/j.spl.2020.108854}.
\href{https://arxiv.org/abs/1908.03428}{arXiv:1908.03428}.

\bibitem{Wang}
M. Wang, \emph{Moments of Gamma type and three-parametric
Mittag--Leffler function}, preprint, 2024.
\href{https://arxiv.org/abs/2410.19330}{arXiv:2410.19330}.

\bibitem{DLMF}
NIST Digital Library of Mathematical Functions,
\emph{Section 5.11: Asymptotic Expansions},
\url{https://dlmf.nist.gov/5.11}. Accessed 3 November 2026.

\bibitem{Corpus}
QEEDDD, \emph{Verifiable Mathematical Reasoning Corpus},
GitHub repository, snapshot of 2 November 2026, commit
\texttt{ecd87de3348dcb706a76bb12fde722d7a348a012}.
\url{https://github.com/QEEDDD/Verifiable-Mathematical-Reasoning-Corpus}.
Methodological comparisons: manuscripts [067] and [096].

\end{thebibliography}
\end{document}
